% GENERATED FILE. Edit canonical lessons, then run npm run generate:books.
% canonical-source-hash: fnv1a64:c5a25383aeb4778a
% canonical-lessons: PA-C329-sahura, PA-C329-sass, PA-C329-numh, PA-C329-javai, PA-C329-kuram, PA-R329-review-323-329

\chapter{Parents-in-Law}
\label{ch:pa-a2-parents-in-law}
\hlchaptermodality{329}

\begin{chapteropening}
\textbf{By the end of this chapter:} \emph{I can name my in-laws and say who they are.}

Complete the last lesson of chapter 329: I can name my in-laws and say who they are.
\end{chapteropening}

\section[\texorpdfstring{\pa{ਸਹੁਰਾ} (sahurā)}{ਸਹੁਰਾ (sahurā)}]{\pa{ਸਹੁਰਾ} (sahurā) — a father-in-law}

\label{lesson:PA-C329-sahura}



\begin{quote}
Before the new one: say the Punjabi for to press gently, then the Punjabi for to put up with.
\end{quote}

\subsection*{You'll want to know: \pa{ਸਹੁਰਾ}}
\textbf{\pa{ਸਹੁਰਾ}} — \emph{sahurā} — \textquotedblleft{}a father-in-law\textquotedblright{}.

\begin{quote}
\pa{ਮੇਰਾ} \pa{ਸਹੁਰਾ} \pa{ਇੱਕ} \pa{ਅਧਿਆਪਕ} \pa{ਹੈ।}
\end{quote}

\begin{quote}
\emph{merā sahurā ikk adhiāpak hai.}
\end{quote}

\begin{quote}
My father-in-law is a teacher.
\end{quote}

A masculine noun.

\begin{cousinweb}
\textbf{\pa{ਸਹੁਰਾ}} (\emph{sahurā}) is inherited from Sanskrit \emph{śvaśura}, \textquotedblleft{}a husband's or wife's father\textquotedblright{}; CDIAL lists Punjabi \emph{sahurā}, \textquotedblleft{}father-in-law\textquotedblright{}, under it (CDIAL 12753).
\end{cousinweb}

\begin{grammarlens}[title={What you've built}]
One more word for this chapter.
\end{grammarlens}

\subsection*{Your turn}
\begin{itemize}
\raggedright
  \item \emph{Say it:} \emph{sahurā}
  \item \emph{Say it:} \emph{sahurā}, once more
  \item \emph{Say it:} \emph{merā sahurā ikk adhiāpak hai.} — My father-in-law is a teacher.
  \item \emph{Recall:} say the Punjabi for to press gently, then the Punjabi for to put up with, then say \emph{sahurā} again
\end{itemize}

\begin{tcolorbox}[breakable,colback=teal!4,colframe=teal!35!black,title={Before you move on}]
What does \pa{ਸਹੁਰਾ} mean? (\textquotedblleft{}A father-in-law\textquotedblright{}.) Say it once more.
\end{tcolorbox}
\section[\texorpdfstring{\pa{ਸੱਸ} (sass)}{ਸੱਸ (sass)}]{\pa{ਸੱਸ} (sass) — a mother-in-law}

\label{lesson:PA-C329-sass}



\begin{quote}
Before the new one: say the Punjabi for to put up with, then the Punjabi for a father-in-law.
\end{quote}

\subsection*{You'll want to know: \pa{ਸੱਸ}}
\textbf{\pa{ਸੱਸ}} — \emph{sass} — \textquotedblleft{}a mother-in-law\textquotedblright{}.

\begin{quote}
\pa{ਮੇਰੀ} \pa{ਸੱਸ} \pa{ਬਹੁਤ} \pa{ਚੰਗਾ} \pa{ਖਾਣਾ} \pa{ਪਕਾਉਂਦੀ} \pa{ਹੈ।}
\end{quote}

\begin{quote}
\emph{merī sass bahut changā khāṇā pakāundī hai.}
\end{quote}

\begin{quote}
My mother-in-law cooks very good food.
\end{quote}

A feminine noun.

\begin{cousinweb}
\textbf{\pa{ਸੱਸ}} (\emph{sass}) is inherited from Sanskrit \emph{śvaśrū}, \textquotedblleft{}a husband's or wife's mother\textquotedblright{}; CDIAL lists Punjabi \emph{sass}, \textquotedblleft{}husband's or wife's mother\textquotedblright{}, under it (CDIAL 12759).
\end{cousinweb}

\begin{grammarlens}[title={What you've built}]
One more word for this chapter.
\end{grammarlens}

\subsection*{Your turn}
\begin{itemize}
\raggedright
  \item \emph{Say it:} \emph{sass}
  \item \emph{Say it:} \emph{sass}, once more
  \item \emph{Say it:} \emph{merī sass bahut changā khāṇā pakāundī hai.} — My mother-in-law cooks very good food.
  \item \emph{Recall:} say the Punjabi for to put up with, then the Punjabi for a father-in-law, then say \emph{sass} again
  \item \emph{Write it:} copy \textbf{\pa{ਸੱਸ}} once, checking the whole word
\end{itemize}

\begin{tcolorbox}[breakable,colback=teal!4,colframe=teal!35!black,title={Before you move on}]
What does \pa{ਸੱਸ} mean? (\textquotedblleft{}A mother-in-law\textquotedblright{}.) Say it once more.
\end{tcolorbox}
\section[\texorpdfstring{\pa{ਨੂੰਹ} (nūṁh)}{ਨੂੰਹ (nūṁh)}]{\pa{ਨੂੰਹ} (nūṁh) — a daughter-in-law}

\label{lesson:PA-C329-numh}



\begin{quote}
Before the new one: say the Punjabi for a father-in-law, then the Punjabi for a mother-in-law.
\end{quote}

\subsection*{You'll want to know: \pa{ਨੂੰਹ}}
\textbf{\pa{ਨੂੰਹ}} — \emph{nūṁh} — \textquotedblleft{}a daughter-in-law\textquotedblright{}.

\begin{quote}
\pa{ਸਾਡੀ} \pa{ਨੂੰਹ} \pa{ਰੋਜ਼} \pa{ਸਵੇਰੇ} \pa{ਸੈਰ} \pa{ਕਰਦੀ} \pa{ਹੈ।}
\end{quote}

\begin{quote}
\emph{sāḍī nūṁh roz savere sair kardī hai.}
\end{quote}

\begin{quote}
Our daughter-in-law goes for a walk every morning.
\end{quote}

A feminine noun.

\begin{cousinweb}
\textbf{\pa{ਨੂੰਹ}} (\emph{nūṁh}) is inherited from Sanskrit \emph{snuṣā}, \textquotedblleft{}a son's wife\textquotedblright{}; CDIAL lists Punjabi \emph{nū̃h}, \textquotedblleft{}son's wife\textquotedblright{}, under it (CDIAL 13801).
\end{cousinweb}

\begin{grammarlens}[title={What you've built}]
One more word for this chapter.
\end{grammarlens}

\subsection*{Your turn}
\begin{itemize}
\raggedright
  \item \emph{Say it:} \emph{nūṁh}
  \item \emph{Say it:} \emph{nūṁh}, once more
  \item \emph{Say it:} \emph{sāḍī nūṁh roz savere sair kardī hai.} — Our daughter-in-law goes for a walk every morning.
  \item \emph{Recall:} say the Punjabi for a father-in-law, then the Punjabi for a mother-in-law, then say \emph{nūṁh} again
  \item \emph{Write it:} copy \textbf{\pa{ਨੂੰਹ}} once, checking the whole word
\end{itemize}

\begin{tcolorbox}[breakable,colback=teal!4,colframe=teal!35!black,title={Before you move on}]
What does \pa{ਨੂੰਹ} mean? (\textquotedblleft{}A daughter-in-law\textquotedblright{}.) Say it once more.
\end{tcolorbox}
\section[\texorpdfstring{\pa{ਜਵਾਈ} (javāī)}{ਜਵਾਈ (javāī)}]{\pa{ਜਵਾਈ} (javāī) — a son-in-law}

\label{lesson:PA-C329-javai}



\begin{quote}
Before the new one: say the Punjabi for a mother-in-law, then the Punjabi for a daughter-in-law.
\end{quote}

\subsection*{You'll want to know: \pa{ਜਵਾਈ}}
\textbf{\pa{ਜਵਾਈ}} — \emph{javāī} — \textquotedblleft{}a son-in-law\textquotedblright{}.

\begin{quote}
\pa{ਸਾਡਾ} \pa{ਜਵਾਈ} \pa{ਬੈਂਕ} \pa{ਵਿੱਚ} \pa{ਕੰਮ} \pa{ਕਰਦਾ} \pa{ਹੈ।}
\end{quote}

\begin{quote}
\emph{sāḍā javāī baiṅk vich kamm kardā hai.}
\end{quote}

\begin{quote}
Our son-in-law works in a bank.
\end{quote}

A masculine noun.

\begin{cousinweb}
\textbf{\pa{ਜਵਾਈ}} (\emph{javāī}) is inherited from Sanskrit \emph{jāmātṛ}, \textquotedblleft{}a daughter's husband\textquotedblright{}; CDIAL lists Punjabi \emph{javāī}, \textquotedblleft{}daughter's husband\textquotedblright{}, under it (CDIAL 5198).
\end{cousinweb}

\begin{grammarlens}[title={What you've built}]
One more word for this chapter.
\end{grammarlens}

\subsection*{Your turn}
\begin{itemize}
\raggedright
  \item \emph{Say it:} \emph{javāī}
  \item \emph{Say it:} \emph{javāī}, once more
  \item \emph{Say it:} \emph{sāḍā javāī baiṅk vich kamm kardā hai.} — Our son-in-law works in a bank.
  \item \emph{Recall:} say the Punjabi for a mother-in-law, then the Punjabi for a daughter-in-law, then say \emph{javāī} again
  \item \emph{Write it:} copy \textbf{\pa{ਜਵਾਈ}} once, checking the whole word
\end{itemize}

\begin{tcolorbox}[breakable,colback=teal!4,colframe=teal!35!black,title={Before you move on}]
What does \pa{ਜਵਾਈ} mean? (\textquotedblleft{}A son-in-law\textquotedblright{}.) Say it once more.
\end{tcolorbox}
\section[\texorpdfstring{\pa{ਕੁੜਮ} (kuṛam)}{ਕੁੜਮ (kuṛam)}]{\pa{ਕੁੜਮ} (kuṛam) — a child's father-in-law}

\label{lesson:PA-C329-kuram}



\begin{quote}
Before the new one: say the Punjabi for a daughter-in-law, then the Punjabi for a son-in-law.
\end{quote}

\subsection*{You'll want to know: \pa{ਕੁੜਮ}}
\textbf{\pa{ਕੁੜਮ}} — \emph{kuṛam} — \textquotedblleft{}a child's father-in-law\textquotedblright{}.

\begin{quote}
\pa{ਸਾਡੇ} \pa{ਕੁੜਮ} \pa{ਕੱਲ੍ਹ} \pa{ਸਾਡੇ} \pa{ਘਰ} \pa{ਆਏ।}
\end{quote}

\begin{quote}
\emph{sāḍe kuṛam kall sāḍe ghar āe.}
\end{quote}

\begin{quote}
Our daughter's father-in-law came to our house yesterday.
\end{quote}

A masculine noun.

\begin{cousinweb}
\textbf{\pa{ਕੁੜਮ}} (\emph{kuṛam}) is inherited from Sanskrit \emph{kuṭumbin}, \textquotedblleft{}a householder\textquotedblright{}; CDIAL lists Punjabi \emph{kuṛam}, \textquotedblleft{}a relation by marriage, child's father-in-law\textquotedblright{}, under it (CDIAL 3235).
\end{cousinweb}

\begin{grammarlens}[title={What you've built}]
One more word for this chapter.
\end{grammarlens}

\subsection*{Your turn}
\begin{itemize}
\raggedright
  \item \emph{Say it:} \emph{kuṛam}
  \item \emph{Say it:} \emph{kuṛam}, once more
  \item \emph{Say it:} \emph{sāḍe kuṛam kall sāḍe ghar āe.} — Our daughter's father-in-law came to our house yesterday.
  \item \emph{Recall:} say the Punjabi for a daughter-in-law, then the Punjabi for a son-in-law, then say \emph{kuṛam} again
  \item \emph{Write it:} copy \textbf{\pa{ਕੁੜਮ}} once, checking the whole word
\end{itemize}

\begin{tcolorbox}[breakable,colback=teal!4,colframe=teal!35!black,title={Before you move on}]
What does \pa{ਕੁੜਮ} mean? (\textquotedblleft{}A child's father-in-law\textquotedblright{}.) Say it once more.
\end{tcolorbox}
\section[(dialogue)]{Looking back: from Inside the Body to Parents-in-Law}

\label{lesson:PA-R329-review-323-329}



\begin{quote}
Say the Punjabi for the belly and a bone.
\end{quote}

\subsection*{Your turn}
Say these aloud:

\begin{itemize}
\raggedright
  \item the belly, a bone, the liver, a rib, to ache
  \item an arm, a shoulder, to pant, sweat, a dizzy spell
  \item to improve, to get worse, keeping off certain food, to chew, an ache
  \item to be afraid, anger, surprised, worry, embarrassment
  \item to sulk, to get annoyed, alone, sorrow, a teardrop
  \item courage, upset, reassurance, to press gently, to put up with
  \item a father-in-law, a mother-in-law, a daughter-in-law, a son-in-law, a child's father-in-law
\end{itemize}

\begin{tcolorbox}[breakable,colback=teal!4,colframe=teal!35!black,title={Before you move on}]
The belly? (\textbf{\pa{ਢਿੱਡ}}, \emph{ḍhiḍḍ}.) A bone? (\textbf{\pa{ਹੱਡੀ}}, \emph{haḍḍī}.) The liver? (\textbf{\pa{ਕਲੇਜਾ}}, \emph{kalejā}.) A rib? (\textbf{\pa{ਪਸਲੀ}}, \emph{paslī}.) To ache? (\textbf{\pa{ਦੁਖਣਾ}}, \emph{dukhṇā}.) An arm? (\textbf{\pa{ਬਾਂਹ}}, \emph{bāṁh}.) A shoulder? (\textbf{\pa{ਮੋਢਾ}}, \emph{moḍhā}.) To pant? (\textbf{\pa{ਹਫਣਾ}}, \emph{haphṇā}.) Sweat? (\textbf{\pa{ਪਸੀਨਾ}}, \emph{pasīnā}.) A dizzy spell? (\textbf{\pa{ਚੱਕਰ}}, \emph{chakkar}.) To improve? (\textbf{\pa{ਸੁਧਰਨਾ}}, \emph{sudharnā}.) To get worse? (\textbf{\pa{ਵਿਗੜਨਾ}}, \emph{vigaṛnā}.) Keeping off certain food? (\textbf{\pa{ਪਰਹੇਜ਼}}, \emph{parhez}.) To chew? (\textbf{\pa{ਚਬਾਉਣਾ}}, \emph{chabāuṇā}.) An ache? (\textbf{\pa{ਪੀੜ}}, \emph{pīṛ}.) To be afraid? (\textbf{\pa{ਡਰਨਾ}}, \emph{ḍarnā}.) Anger? (\textbf{\pa{ਗੁੱਸਾ}}, \emph{gussā}.) Surprised? (\textbf{\pa{ਹੈਰਾਨ}}, \emph{hairān}.) Worry? (\textbf{\pa{ਫ਼ਿਕਰ}}, \emph{fikar}.) Embarrassment? (\textbf{\pa{ਸ਼ਰਮ}}, \emph{sharam}.) To sulk? (\textbf{\pa{ਰੁੱਸਣਾ}}, \emph{russṇā}.) To get annoyed? (\textbf{\pa{ਖਿਝਣਾ}}, \emph{khijhṇā}.) Alone? (\textbf{\pa{ਇਕੱਲਾ}}, \emph{ikallā}.) Sorrow? (\textbf{\pa{ਦੁੱਖ}}, \emph{dukkh}.) A teardrop? (\textbf{\pa{ਹੰਝੂ}}, \emph{hanjhū}.) Courage? (\textbf{\pa{ਹੌਸਲਾ}}, \emph{hauslā}.) Upset? (\textbf{\pa{ਪਰੇਸ਼ਾਨ}}, \emph{pareshān}.) Reassurance? (\textbf{\pa{ਤਸੱਲੀ}}, \emph{tasallī}.) To press gently? (\textbf{\pa{ਦਬਾਉਣਾ}}, \emph{dabāuṇā}.) To put up with? (\textbf{\pa{ਸਹਿਣਾ}}, \emph{sahiṇā}.) A father-in-law? (\textbf{\pa{ਸਹੁਰਾ}}, \emph{sahurā}.) A mother-in-law? (\textbf{\pa{ਸੱਸ}}, \emph{sass}.) A daughter-in-law? (\textbf{\pa{ਨੂੰਹ}}, \emph{nūṁh}.) A son-in-law? (\textbf{\pa{ਜਵਾਈ}}, \emph{javāī}.) A child's father-in-law? (\textbf{\pa{ਕੁੜਮ}}, \emph{kuṛam}.)
\end{tcolorbox}
